% GENERATED FILE. Edit canonical lessons, then run npm run generate:books.
% canonical-source-hash: fnv1a64:fc0ad3f3b87b6db1
% canonical-lessons: AR-C428-sikkat-hadid, AR-C428-araba, AR-C428-maqsura, AR-C428-mitru, AR-C428-anzala

\chapter{On the railway}
\label{ch:ar-on-the-railway}
\hlchaptermodality{428}

\begin{chapteropening}
\textbf{By the end of this chapter:} \emph{I can say a railway, a carriage, a compartment, the metro and to unload.}

Complete the last lesson of chapter 428: I can say a railway, a carriage, a compartment, the metro and to unload.
\end{chapteropening}

\section[\texorpdfstring{\ar{سكة} \ar{حديد} (sikkat ḥadīd)}{سكة حديد (sikkat ḥadīd)}]{\ar{سكة} \ar{حديد} (sikkat ḥadīd) — a railway}

\label{lesson:AR-C428-sikkat-hadid}



\begin{quote}
Before the new one: say the Arabic for to hire, then the Arabic for a backpack.
\end{quote}

\subsection*{You'll want to know: \ar{سكة} \ar{حديد}}
\textbf{\ar{سكة} \ar{حديد}} — \emph{sikkat ḥadīd} — \textquotedblleft{}a railway\textquotedblright{}.

\begin{quote}
\ar{وصلت} \ar{سكة} \ar{الحديد} \ar{إلى} \ar{قريتنا} \ar{قبل} \ar{قرن}.
\end{quote}

\begin{quote}
\emph{waṣalat sikkat al-ḥadīd ilā qaryatinā qabla qarn.}
\end{quote}

\begin{quote}
The railway reached our village a century ago.
\end{quote}

\begin{cousinweb}
\textbf{\ar{سكة} \ar{حديد}} (\emph{sikkat ḥadīd}) is built from two native Arabic words: \emph{sikka}, a straight road or row, on the root \emph{s-k-k} (Lane), and \emph{ḥadīd}, \textquotedblleft{}iron\textquotedblright{}, on the root \emph{ḥ-d-d} (Lane): a road of iron.
\end{cousinweb}

\begin{grammarlens}[title={What you've built}]
One more word for getting around.
\end{grammarlens}

\subsection*{Your turn}
\begin{itemize}
\raggedright
  \item \emph{Say it:} \emph{sikkat ḥadīd}
  \item \emph{Say it:} \emph{sikkat ḥadīd}, once more
  \item \emph{Say it:} \emph{waṣalat sikkat al-ḥadīd ilā qaryatinā qabla qarn.} — The railway reached our village a century ago.
  \item \emph{Recall:} say the Arabic for to hire, then the Arabic for a backpack, then say \emph{sikkat ḥadīd} again
\end{itemize}

\begin{tcolorbox}[breakable,colback=teal!4,colframe=teal!35!black,title={Before you move on}]
What does \ar{سكة} \ar{حديد} mean? (\textquotedblleft{}A railway\textquotedblright{}.) Say it once more.
\end{tcolorbox}
\section[\texorpdfstring{\ar{عربة} (ʿaraba)}{عربة (ʿaraba)}]{\ar{عربة} (ʿaraba) — a carriage; a cart}

\label{lesson:AR-C428-araba}



\begin{quote}
Before the new one: say the Arabic for a backpack, then the Arabic for a railway.
\end{quote}

\subsection*{You'll want to know: \ar{عربة}}
\textbf{\ar{عربة}} — \emph{ʿaraba} — \textquotedblleft{}a carriage; a cart\textquotedblright{}.

\begin{quote}
\ar{جلسنا} \ar{في} \ar{العربة} \ar{الأخيرة} \ar{من} \ar{القطار}.
\end{quote}

\begin{quote}
\emph{jalasnā fī al-ʿaraba al-akhīra min al-qiṭār.}
\end{quote}

\begin{quote}
We sat in the last carriage of the train.
\end{quote}

\begin{cousinweb}
\textbf{\ar{عربة}} (\emph{ʿaraba}) names a cart and a railway carriage; the origin is not settled, and some trace it to Turkish \emph{araba}.
\end{cousinweb}

\begin{grammarlens}[title={What you've built}]
One more word for getting around.
\end{grammarlens}

\subsection*{Your turn}
\begin{itemize}
\raggedright
  \item \emph{Say it:} \emph{ʿaraba}
  \item \emph{Say it:} \emph{ʿaraba}, once more
  \item \emph{Say it:} \emph{jalasnā fī al-ʿaraba al-akhīra min al-qiṭār.} — We sat in the last carriage of the train.
  \item \emph{Recall:} say the Arabic for a backpack, then the Arabic for a railway, then say \emph{ʿaraba} again
\end{itemize}

\begin{tcolorbox}[breakable,colback=teal!4,colframe=teal!35!black,title={Before you move on}]
What does \ar{عربة} mean? (\textquotedblleft{}A carriage; a cart\textquotedblright{}.) Say it once more.
\end{tcolorbox}
\section[\texorpdfstring{\ar{مقصورة} (maqṣūra)}{مقصورة (maqṣūra)}]{\ar{مقصورة} (maqṣūra) — a compartment; a cabin}

\label{lesson:AR-C428-maqsura}



\begin{quote}
Before the new one: say the Arabic for a railway, then the Arabic for a carriage.
\end{quote}

\subsection*{You'll want to know: \ar{مقصورة}}
\textbf{\ar{مقصورة}} — \emph{maqṣūra} — \textquotedblleft{}a compartment; a cabin\textquotedblright{}.

\begin{quote}
\ar{كانت} \ar{المقصورة} \ar{مزدحمة} \ar{بالمسافرين}.
\end{quote}

\begin{quote}
\emph{kānat al-maqṣūra muzdaḥima bi-l-musāfirīn.}
\end{quote}

\begin{quote}
The compartment was crowded with travellers.
\end{quote}

\begin{cousinweb}
\textbf{\ar{مقصورة}} (\emph{maqṣūra}) is built on the native Arabic root \emph{q-ṣ-r}, \textquotedblleft{}to confine, to keep within bounds\textquotedblright{} (Lane). It is a passive participle: a space closed off, a compartment.
\end{cousinweb}

\begin{grammarlens}[title={What you've built}]
One more word for getting around.
\end{grammarlens}

\subsection*{Your turn}
\begin{itemize}
\raggedright
  \item \emph{Say it:} \emph{maqṣūra}
  \item \emph{Say it:} \emph{maqṣūra}, once more
  \item \emph{Say it:} \emph{kānat al-maqṣūra muzdaḥima bi-l-musāfirīn.} — The compartment was crowded with travellers.
  \item \emph{Recall:} say the Arabic for a railway, then the Arabic for a carriage, then say \emph{maqṣūra} again
\end{itemize}

\begin{tcolorbox}[breakable,colback=teal!4,colframe=teal!35!black,title={Before you move on}]
What does \ar{مقصورة} mean? (\textquotedblleft{}A compartment; a cabin\textquotedblright{}.) Say it once more.
\end{tcolorbox}
\section[\texorpdfstring{\ar{مترو} (mitrū)}{مترو (mitrū)}]{\ar{مترو} (mitrū) — the metro; the underground}

\label{lesson:AR-C428-mitru}



\begin{quote}
Before the new one: say the Arabic for a carriage, then the Arabic for a compartment.
\end{quote}

\subsection*{You'll want to know: \ar{مترو}}
\textbf{\ar{مترو}} — \emph{mitrū} — \textquotedblleft{}the metro; the underground\textquotedblright{}.

\begin{quote}
\ar{يذهب} \ar{أبي} \ar{إلى} \ar{العمل} \ar{بالمترو}.
\end{quote}

\begin{quote}
\emph{yadhhabu abī ilā al-ʿamal bi-l-mitrū.}
\end{quote}

\begin{quote}
My father goes to work by metro.
\end{quote}

\begin{cousinweb}
\textbf{\ar{مترو}} (\emph{mitrū}) is usually traced to French \emph{métro}, short for \emph{métropolitain}.
\end{cousinweb}

\begin{grammarlens}[title={What you've built}]
One more word for getting around.
\end{grammarlens}

\subsection*{Your turn}
\begin{itemize}
\raggedright
  \item \emph{Say it:} \emph{mitrū}
  \item \emph{Say it:} \emph{mitrū}, once more
  \item \emph{Say it:} \emph{yadhhabu abī ilā al-ʿamal bi-l-mitrū.} — My father goes to work by metro.
  \item \emph{Recall:} say the Arabic for a carriage, then the Arabic for a compartment, then say \emph{mitrū} again
\end{itemize}

\begin{tcolorbox}[breakable,colback=teal!4,colframe=teal!35!black,title={Before you move on}]
What does \ar{مترو} mean? (\textquotedblleft{}The metro; the underground\textquotedblright{}.) Say it once more.
\end{tcolorbox}
\section[\texorpdfstring{\ar{أنزل} (anzala)}{أنزل (anzala)}]{\ar{أنزل} (anzala) — to unload; to bring down}

\label{lesson:AR-C428-anzala}



\begin{quote}
Before the new one: say the Arabic for a compartment, then the Arabic for the metro.
\end{quote}

\subsection*{You'll want to know: \ar{أنزل}}
\textbf{\ar{أنزل}} — \emph{anzala} — \textquotedblleft{}to unload; to bring down\textquotedblright{}.

\begin{quote}
\ar{أنزل} \ar{أبي} \ar{الأمتعة} \ar{من} \ar{القطار}.
\end{quote}

\begin{quote}
\emph{anzala abī al-amtiʿa min al-qiṭār.}
\end{quote}

\begin{quote}
My father unloaded the luggage from the train.
\end{quote}

Present: \ar{ينزل}, \emph{yunzil}.

Unvocalized, its present \ar{ينزل} looks like that of \emph{nazala}, \textquotedblleft{}to go down\textquotedblright{} (\emph{yanzil}).

\begin{cousinweb}
\textbf{\ar{أنزل}} (\emph{anzala}) is built on the native Arabic root \emph{n-z-l}, \textquotedblleft{}to come down, to alight\textquotedblright{} (Lane). It is Form IV, \textquotedblleft{}to bring down\textquotedblright{}.
\end{cousinweb}

\begin{grammarlens}[title={What you've built}]
One more word for getting around.
\end{grammarlens}

\subsection*{Your turn}
\begin{itemize}
\raggedright
  \item \emph{Say it:} \emph{anzala}
  \item \emph{Say it:} \emph{anzala}, once more
  \item \emph{Say it:} \emph{anzala abī al-amtiʿa min al-qiṭār.} — My father unloaded the luggage from the train.
  \item \emph{Recall:} say the Arabic for a compartment, then the Arabic for the metro, then say \emph{anzala} again
\end{itemize}

\begin{tcolorbox}[breakable,colback=teal!4,colframe=teal!35!black,title={Before you move on}]
What does \ar{أنزل} mean? (\textquotedblleft{}To unload; to bring down\textquotedblright{}.) Say it once more.
\end{tcolorbox}
